\documentclass[11pt]{article}

\usepackage[margin=1in]{geometry}
\usepackage{amsmath,amssymb,amsthm}
\usepackage{enumitem}
\usepackage{hyperref}
\usepackage{xcolor}
\usepackage{booktabs}

\hypersetup{colorlinks=true, linkcolor=blue, urlcolor=blue, citecolor=blue}

\title{TOAE209/TOAH209 -- Design and Analysis of Algorithms\\[4pt]
\large Week 4: Theoretical Questions}
\author{Foreign Trade University}
\date{}

\begin{document}
\maketitle

\section*{Instructions}

Part A contains 22 questions requiring written explanations or short proofs. Part B is a
10-question quiz (true/false and multiple choice). Attempt every question on your own
before consulting Part C (Solutions) at the end of this document.

\section*{Part A: Theory Questions}

\begin{enumerate}[label=\textbf{A\arabic*.}, leftmargin=2.2em]

\item What is a \emph{greedy algorithm}? Describe, in general terms, the structure of a
greedy algorithm: what choice does it make at each step, and what does it never do
afterward?

\item Greedy algorithms are attractive because they are usually simple and efficient, but
this efficiency comes with a risk. Explain what can go wrong if a greedy rule is applied
without a correctness proof, and why ``it worked on my test cases'' is not sufficient
justification for a greedy algorithm.

\item Describe the \emph{``greedy stays ahead''} proof technique in general terms: what
quantity do we compare between the greedy solution and an arbitrary optimal solution, and
what do we conclude by induction?

\item Describe the \emph{exchange argument} proof technique in general terms: starting
from an arbitrary optimal solution $O$, what transformation is applied to $O$, and what
must be shown about the effect of this transformation on $O$'s objective value?

\item State the Interval Scheduling Problem precisely: what is the input (in terms of
start and finish times), what is a \emph{feasible} (compatible) set of requests, and what
is the objective?

\item List three natural greedy heuristics for Interval Scheduling that \emph{do not}
always produce an optimal solution (e.g.\ based on start time, length, or number of
conflicts), and for each, briefly describe the kind of instance that defeats it.

\item State the Earliest-Finish-Time-First (EFT) algorithm for Interval Scheduling.
Why is sorting by \emph{finish} time (rather than start time) the key idea that makes
this greedy rule work?

\item Prove that the Earliest-Finish-Time-First algorithm produces an optimal solution to
Interval Scheduling, using the ``greedy stays ahead'' technique. State precisely what is
being compared at each step of the induction.

\item State the Interval Partitioning (Interval Graph Coloring) Problem precisely: what is
the input, what is a \emph{resource assignment}, what makes an assignment \emph{valid},
and what is the objective?

\item Define the \emph{depth} of a set of intervals. Prove that no valid resource
assignment can use fewer than $\mathrm{depth}$ resources (i.e.\ depth is a lower bound on
the optimum).

\item State the Earliest-Start-Time-First algorithm for Interval Partitioning. Prove that
it never needs more than $\mathrm{depth}$ resources, and conclude (using Q10) that it is
optimal.

\item State the Scheduling to Minimize Lateness Problem precisely: what is the input (in
terms of processing times and deadlines), what is a \emph{schedule}, how is the
\emph{lateness} of a job defined, and what is the objective?

\item State the Earliest-Deadline-First (EDF) algorithm for Minimize Lateness. Give an
example showing that Shortest-Job-First (sorting by processing time $t_i$ instead of
deadline $d_i$) can produce a strictly larger maximum lateness than EDF.

\item Define an \emph{inversion} in a schedule (with respect to deadlines) and define
\emph{idle time}. Prove that an optimal schedule for Minimize Lateness can always be
chosen to have no idle time and no inversions.

\item Using Q14, complete the proof that Earliest-Deadline-First is optimal for Minimize
Lateness. Your argument should explain why EDF itself has no inversions and no idle time,
and why any schedule with no inversions and no idle time has the same maximum lateness as
any other such schedule (up to the specific job achieving the maximum).

\item State the Optimal (Offline) Caching Problem precisely: what is the input (a sequence
of page requests and a cache of size $k$), what is a \emph{cache miss}, and what is the
objective?

\item State Belady's Furthest-In-Future (FIF) algorithm. Explain, at a high level, why
evicting the page whose next use is furthest in the future (or that is never used again)
is the ``safe'' eviction choice -- i.e.\ why an exchange argument can convert any optimal
schedule into one that makes the same eviction at this step without increasing the number
of misses.

\item Define the Least-Recently-Used (LRU) eviction policy. Explain why LRU is a
reasonable \emph{online} heuristic (no knowledge of the future) even though it is not
always optimal, and give a short example sequence where LRU and FIF make different
eviction decisions.

\item The Coin Change problem asks for the minimum number of coins (with unlimited supply
of each denomination) summing to a target amount $V$. The greedy algorithm repeatedly
takes the largest denomination $\le$ the remaining amount. Define what it means for a
coin system to be \emph{canonical} (greedy-optimal for all amounts), and explain why the
standard system $\{1, 5, 10, 25\}$ (US cents) is canonical while $\{1, 3, 4\}$ is not
(give the smallest counter-example amount for $\{1,3,4\}$).

\item State the Job Sequencing with Deadlines and Profits problem: each job $i$ has a
deadline $d_i$ and profit $p_i$, all jobs take one unit of time, and a job earns its
profit only if scheduled at or before its deadline on a single machine. Describe the
greedy algorithm (sort by profit, then try to place each job as late as possible before
its deadline) and explain at a high level why processing jobs in decreasing order of
profit and placing each as late as possible is the right greedy rule.

\item Compare the four ``flavors'' of greedy correctness seen this week: (a)
Interval Scheduling (EFT, exchange/stays-ahead), (b) Interval Partitioning (depth lower
bound + matching upper bound), (c) Minimize Lateness (exchange argument removing
inversions), and (d) Optimal Caching (exchange argument on eviction choice). What do all
four proofs have in common, structurally?

\item \emph{Discussion (open-ended).} Suppose you are given a new optimization problem
and you come up with a greedy rule that seems to work on every small example you try by
hand. Describe a systematic process (drawing on the techniques in this chapter) you would
follow to either (a) gain confidence that the rule is correct and find a proof, or (b)
discover a counter-example that shows it is wrong. Be specific about what kinds of
instances you would try to construct.

\end{enumerate}

\newpage
\section*{Part B: Quiz}

\begin{enumerate}[label=\textbf{B\arabic*.}, leftmargin=2.2em]

\item \textbf{(True/False)} A greedy algorithm, once it makes a choice, may later revisit
and revise that choice if a better option becomes available.

\item \textbf{(Multiple choice)} The Earliest-Finish-Time-First algorithm for Interval
Scheduling sorts requests by:
\begin{enumerate}[label=(\alph*)]
    \item Start time, increasing
    \item Finish time, increasing
    \item Length (duration), increasing
    \item Number of conflicts, increasing
\end{enumerate}

\item \textbf{(Short answer)} Consider three intervals: $[0,10]$, $[1,2]$, and $[3,4]$.
Shortest-interval-first would pick $[1,2]$ first (and then $[3,4]$, both compatible with
$[1,2]$ but not $[0,10]$). What is the size of this shortest-first solution, versus the
true optimum?

\item \textbf{(True/False)} The depth of a set of intervals is always a valid lower bound
on the number of resources (e.g.\ classrooms) needed for interval partitioning, regardless
of which algorithm is used to construct the assignment.

\item \textbf{(Multiple choice)} In Scheduling to Minimize Lateness, the Earliest-Deadline-
First algorithm sorts jobs by:
\begin{enumerate}[label=(\alph*)]
    \item Processing time $t_i$, increasing
    \item Deadline $d_i$, increasing
    \item Profit, decreasing
    \item Slack $d_i - t_i$, decreasing
\end{enumerate}

\item \textbf{(True/False)} An optimal schedule for Minimize Lateness can always be
transformed (without increasing the maximum lateness) into one with no idle time and no
inversions.

\item \textbf{(Short answer)} In Belady's Furthest-In-Future caching algorithm, when a
cache miss occurs and the cache is full, which page is evicted?

\item \textbf{(Multiple choice)} Least-Recently-Used (LRU) caching is best described as:
\begin{enumerate}[label=(\alph*)]
    \item An offline algorithm that requires knowledge of all future requests
    \item An online algorithm that uses only the request history seen so far
    \item Always exactly as good as Furthest-In-Future
    \item Optimal only when $k=1$
\end{enumerate}

\item \textbf{(Short answer)} For the coin system $\{1, 3, 4\}$, what is the smallest
amount for which the greedy algorithm (always take the largest coin $\le$ remaining
amount) fails to produce the minimum number of coins, and how many coins does greedy use
vs.\ the optimum?

\item \textbf{(True/False)} In Job Sequencing with Deadlines and Profits (unit-time jobs,
single machine), the greedy algorithm that processes jobs in decreasing order of profit
and schedules each job at the latest available free slot at or before its deadline (if
any) always achieves the maximum total profit.

\end{enumerate}

\newpage
\section*{Part C: Solutions}

\subsection*{Part A}

\begin{enumerate}[label=\textbf{A\arabic*.}, leftmargin=2.2em]

\item A \textbf{greedy algorithm} builds up a solution incrementally by repeatedly making
the choice that looks best \emph{according to some simple, local criterion} at the current
step (e.g.\ ``pick the request that finishes earliest among those still available''). Once
a choice is made, the algorithm \textbf{never reconsiders or undoes it} -- there is no
backtracking. This makes greedy algorithms typically very efficient (often $O(n \log n)$
after sorting), but the local criterion does not automatically guarantee a globally
optimal result; that requires a separate correctness proof.

\item If the greedy rule is wrong, the algorithm may silently produce a \emph{suboptimal}
(or even badly suboptimal) answer on inputs that were not tested -- with no error or
warning, since the algorithm runs to completion and outputs \emph{some} feasible answer.
Testing on a finite set of examples can never rule out the existence of a bad instance
that the test set did not happen to include (this week's exercises construct exactly such
counter-examples for several plausible-looking greedy rules, e.g.\ shortest-interval-first
for Interval Scheduling and shortest-job-first for Minimize Lateness). Only a proof --
typically via exchange argument or greedy-stays-ahead -- rules out \emph{all} possible
inputs simultaneously.

\item In \textbf{greedy stays ahead}, we compare the greedy solution and an arbitrary
optimal solution \emph{step by step} (e.g.\ the $i$-th request selected by greedy vs.\ the
$i$-th request selected by the optimal solution, both sorted in the same order). We prove
by induction on $i$ that, in some measurable sense, greedy's $i$-th choice is ``at least
as good'' as the optimal solution's $i$-th choice (e.g.\ greedy's $i$-th interval finishes
no later). This implies greedy can always be extended at least as far as the optimal
solution, so greedy's final solution is at least as large/good as optimal -- hence
optimal.

\item In the \textbf{exchange argument}, we start from an arbitrary optimal solution $O$
and show that if $O$ differs from the greedy solution $G$ in some specific, local way
(e.g.\ $O$ contains an ``inversion'' that $G$ would never create), we can perform a small
local \emph{exchange} (e.g.\ swap the order of two adjacent jobs, or swap which item is
evicted) that transforms $O$ into another feasible solution $O'$ with objective value
\textbf{no worse} than $O$'s. Repeating this process eventually transforms $O$ into a
solution that looks like (or is exactly as good as) $G$, without ever making the objective
worse -- so $G$ must also be optimal.

\item \textbf{Input:} a set of $n$ requests, each request $i$ specified by a start time
$s_i$ and finish time $f_i$ (with $s_i < f_i$), all referring to a single shared resource.
A set of requests is \textbf{compatible} (feasible) if no two of them overlap in time --
i.e.\ for any two requests $i,j$ in the set, either $f_i \le s_j$ or $f_j \le s_i$. The
\textbf{objective} is to select a compatible set of requests of \textbf{maximum size}
(maximum number of requests accepted).

\item Three plausible but incorrect heuristics:
\begin{itemize}
    \item \textbf{Earliest start time first.} Counter-example: one very long interval
    starting at time 0 (e.g.\ $[0,100]$) blocks out many short intervals that start
    slightly later but would together allow far more requests to be scheduled.
    \item \textbf{Shortest interval (duration) first.} Counter-example: a long interval
    $[0,10]$ and two short intervals $[1,2]$ and $[8,9]$ (mutually compatible, both
    compatible with neither overlapping each other but each overlapping $[0,10]$
    differently) -- e.g.\ $[0,10]$, $[1,2]$, $[3,4]$: shortest-first picks $[1,2]$ then
    $[3,4]$ (size 2, happens to tie optimum here), but with $[0,10]$, $[4,5]$, and three
    short overlapping intervals carefully placed, shortest-first can be made to pick fewer
    than the optimum (this is explored concretely in the exercises).
    \item \textbf{Fewest conflicts first.} Counter-example: an interval with few apparent
    conflicts might still ``use up'' a time slot that many other mutually-compatible short
    intervals need, leading to a smaller final selection than choosing based on finish
    time.
\end{itemize}

\item \textbf{Earliest-Finish-Time-First (EFT):} sort all requests by finish time $f_i$
in increasing order. Initialize the selected set $A = \emptyset$ and a variable
\texttt{last\_finish} $= -\infty$. For each request $i$ in sorted order, if $s_i \ge
\texttt{last\_finish}$ (i.e.\ it starts after the last selected request finishes), add $i$
to $A$ and set \texttt{last\_finish} $= f_i$. Sorting by finish time is the key idea
because it greedily \textbf{preserves the maximum possible remaining time} for future
choices: among all requests that could be selected first, the one that finishes earliest
leaves the largest ``free'' interval $[f_i, \infty)$ available for subsequent selections,
so it can never be strictly worse than any other first choice.

\item \textit{Setup.} Let $G = \{g_1, g_2, \ldots, g_k\}$ be the requests selected by EFT,
ordered by finish time (so $f(g_1) < f(g_2) < \cdots$), and let $O = \{o_1, o_2, \ldots,
o_m\}$ be an arbitrary optimal compatible set, also ordered by finish time. We claim by
induction that for every $i \le \min(k,m)$, $f(g_i) \le f(o_i)$.

\textit{Base case ($i=1$):} $g_1$ is the globally earliest-finishing request among
\emph{all} requests (EFT picks it first, unconditionally), so $f(g_1) \le f(o_1)$.

\textit{Inductive step:} suppose $f(g_i) \le f(o_i)$. Since $o_{i+1}$ is compatible with
$o_i$ (both are in $O$), $s(o_{i+1}) \ge f(o_i) \ge f(g_i)$. So $o_{i+1}$ is a candidate
that EFT could select after $g_i$ (it starts no earlier than $g_i$ finishes). EFT selects
$g_{i+1}$ as the earliest-finishing request compatible with $g_i$ among all remaining
requests, and $o_{i+1}$ is one such candidate, so $f(g_{i+1}) \le f(o_{i+1})$.

\textit{Conclusion.} Suppose for contradiction $k < m$. Then by the claim, $f(g_k) \le
f(o_k)$, and $o_{k+1}$ exists with $s(o_{k+1}) \ge f(o_k) \ge f(g_k)$ -- so $o_{k+1}$ is
compatible with $g_k$ and EFT had not yet terminated (a compatible request remained), a
contradiction (EFT only stops when no compatible request remains). Hence $k \ge m$, and
since $O$ is optimal ($m$ is the maximum possible), $k = m$: EFT is optimal.

\item \textbf{Input:} a set of $n$ requests (intervals), each with a start time $s_i$ and
finish time $f_i$, to be scheduled using some number of identical resources (e.g.\
classrooms, machines). A \textbf{resource assignment} is a function mapping each request
to one of the resources. The assignment is \textbf{valid} if no two requests assigned to
the same resource overlap in time. The \textbf{objective} is to find a valid assignment
using the \textbf{minimum number of resources}.

\item The \textbf{depth} of a set of intervals is the maximum, over all time points $t$,
of the number of intervals that contain $t$ (i.e.\ the maximum number of intervals that
are simultaneously ``active''). \textit{Proof that depth is a lower bound:} let $d$ be the
depth, achieved at some time $t^*$ by a set of $d$ intervals all containing $t^*$. Any two
of these $d$ intervals overlap (both contain $t^*$), so they \emph{must} be assigned to
\emph{different} resources in any valid assignment. Hence any valid assignment needs at
least $d$ distinct resources.

\item \textbf{Earliest-Start-Time-First (for partitioning):} sort all requests by start
time $s_i$ in increasing order. Process requests in this order; for each request,
assign it to any resource that is currently free (i.e.\ not occupied by an
overlapping previously-assigned request) -- if no resource yet used is free, open a new
resource (label) for this request.

\textit{Proof it never needs more than $\mathrm{depth}$ resources:} suppose, for
contradiction, that processing request $i$ in sorted order requires opening resource
number $d+1$ (where $d = \mathrm{depth}$), i.e.\ all $d$ resources opened so far are
currently occupied by intervals overlapping $i$. Each of those $d$ occupying intervals
$j$ has $s_j \le s_i$ (processed earlier, by sorted order) and $f_j > s_i$ (overlaps $i$,
since it is still ``occupying'' its resource at time $s_i$). So all $d$ of those
intervals, plus $i$ itself, contain the time point $s_i$ -- giving $d+1$ intervals all
containing $s_i$, i.e.\ depth $\ge d+1$, contradicting $\mathrm{depth} = d$. So the
algorithm never opens more than $d = \mathrm{depth}$ resources.

\textit{Conclusion:} the algorithm uses at most $\mathrm{depth}$ resources (just proved),
and by Q10 no valid assignment can use fewer than $\mathrm{depth}$ resources. Hence the
algorithm's assignment, using exactly $\mathrm{depth}$ resources (or possibly fewer if
some intervals never co-occur, but never more), is optimal.

\item \textbf{Input:} a set of $n$ jobs to be processed on a single machine (one at a
time, in some order), where job $i$ has a processing time $t_i$ and a deadline $d_i$. A
\textbf{schedule} is an ordering (permutation) of the jobs, run back-to-back starting at
time 0; the \textbf{completion time} $f_i$ of job $i$ is the sum of processing times of
all jobs up to and including $i$ in the order. The \textbf{lateness} of job $i$ is
$\ell_i = \max(0, f_i - d_i)$. The \textbf{objective} is to find a schedule minimizing the
\textbf{maximum lateness} $L = \max_i \ell_i$.

\item \textbf{Earliest-Deadline-First (EDF):} sort jobs by deadline $d_i$ in increasing
order, and schedule them back-to-back in that order (ties broken arbitrarily).

\textit{Counter-example for Shortest-Job-First (SJF):} consider two jobs, job 1 with
$(t_1, d_1) = (1, 100)$ and job 2 with $(t_2, d_2) = (10, 10)$. SJF schedules job 1 first
(shorter processing time), then job 2: job 1 completes at $f_1=1$ (lateness $0$, since
$d_1=100$), job 2 completes at $f_2 = 11$ (lateness $11-10=1$), so $L_{\text{SJF}} = 1$.
EDF schedules job 2 first ($d_2=10 < d_1=100$), then job 1: job 2 completes at $f_2=10$
(lateness $0$), job 1 completes at $f_1=11$ (lateness $11-100 < 0 \Rightarrow 0$), so
$L_{\text{EDF}} = 0 < L_{\text{SJF}} = 1$. SJF is strictly worse here.

\item An \textbf{inversion} is a pair of jobs $i, j$ such that $i$ is scheduled before
$j$ but $d_i > d_j$ (i.e.\ a job with a \emph{later} deadline is scheduled \emph{before} a
job with an \emph{earlier} deadline). \textbf{Idle time} occurs if the machine is left
unused at some point before all jobs are completed, even though jobs remain to be
processed (i.e.\ there is a gap in the schedule before time $\sum_i t_i$).

\textit{No idle time:} removing any idle gap by shifting all subsequent jobs earlier can
only \emph{decrease} (or keep equal) every completion time $f_i$, hence can only decrease
or keep equal every lateness $\ell_i = \max(0, f_i-d_i)$, hence can only decrease or keep
equal $L = \max_i \ell_i$. So an optimal schedule with no idle time always exists.

\textit{No inversions (exchange argument):} suppose an optimal schedule $S$ has an
inversion -- some pair of \emph{adjacent} jobs $i$ (scheduled first) and $j$ (scheduled
right after $i$) with $d_i > d_j$. Swap $i$ and $j$ to get $S'$. All other jobs' completion
times are unchanged. Job $j$ now finishes earlier (where $i$ used to finish), so
$\ell_j$ can only decrease. Job $i$ now finishes at the time $j$ used to finish (later
than before), but since $d_i > d_j$ and $j$'s new completion time equals $i$'s \emph{old}
completion time which was $\le$ ... -- concretely: let $f$ be the time both jobs together
finish (same in $S$ and $S'$). In $S$, $i$ finishes at $f - t_j$ and $j$ at $f$; in $S'$,
$j$ finishes at $f-t_i$ and $i$ at $f$. The lateness of the later-finishing job in each
case is $\max(0, f-d_j)$ for $S$ (job $j$ finishes at $f$, $d_i>d_j$) vs.\ $\max(0,f-d_i)$
for $S'$ (job $i$ finishes at $f$). Since $d_i > d_j$, $f - d_i < f - d_j$, so $S'$'s
``later-finishing-job'' lateness is $\le$ $S$'s. The earlier-finishing job's lateness in
each case is no larger than this. Hence $\max(\ell_i,\ell_j)$ does not increase under the
swap, so the overall maximum lateness $L$ does not increase. Repeatedly removing
inversions this way (each swap strictly reduces the number of inversions, since it fixes
an adjacent inverted pair without creating new inversions elsewhere) terminates with an
inversion-free, idle-time-free schedule that is at least as good as $S$ -- hence also
optimal.

\item EDF, by construction, sorts jobs by increasing deadline and runs them back-to-back,
so it has \textbf{no idle time} (jobs run consecutively from time 0) and \textbf{no
inversions} (jobs are literally ordered by non-decreasing $d_i$, so no job with a larger
deadline precedes one with a smaller deadline). By Q14, \emph{some} optimal schedule has
no idle time and no inversions. Any two schedules with no idle time and no inversions
must order jobs by non-decreasing deadline (ties may be broken arbitrarily among jobs with
equal deadlines, but this does not change any completion times since they all have the
same total processing time up to that point, hence the same lateness values) -- so they
are essentially the same schedule (up to tie-breaking) and achieve the same maximum
lateness. Since EDF is such a schedule and an optimal inversion-free, idle-free schedule
exists with the same maximum lateness, EDF achieves the optimal maximum lateness $L$.

\item \textbf{Input:} a sequence of page requests $r_1, r_2, \ldots, r_n$ (each $r_t$ is a
page identifier) and a cache of fixed size $k$ (can hold $k$ distinct pages at a time). A
\textbf{cache miss} occurs at time $t$ if $r_t$ is not currently in the cache (in which
case $r_t$ must be loaded into the cache, evicting some other page if the cache is full).
The \textbf{objective} is to choose, for each miss, which page to evict, so as to
\textbf{minimize the total number of misses} over the whole sequence.

\item \textbf{Furthest-In-Future (FIF / Belady's algorithm):} on a cache miss with a full
cache, evict the page in the cache whose \emph{next} occurrence in the remaining request
sequence is \emph{furthest in the future} (or evict any page that does not occur again at
all, if such a page is in the cache).

\textit{Why this is ``safe'' (high-level exchange argument):} take any optimal schedule
$O$. If at this step $O$ evicts a different page than FIF would, consider the page $f$
that FIF would evict (whose next use is furthest away) and the page $g$ that $O$ evicts
instead (whose next use is sooner). We can modify $O$ into $O'$ that evicts $f$ instead of
$g$ at this step, and ``patches up'' the schedule afterward (essentially, $O'$ keeps $g$
in cache a bit longer and evicts it later, exactly when $O$ would have needed to bring it
back). Because $f$'s next use is \emph{later} than $g$'s, this patch can never cause
\emph{more} misses than $O$ had -- $O'$ is feasible and has at most as many misses as $O$.
Repeating this argument at every step where FIF and $O$ disagree converts $O$ into a
schedule that behaves like FIF, without ever increasing the miss count -- so FIF is
optimal.

\item \textbf{Least-Recently-Used (LRU):} on a cache miss with a full cache, evict the
page in the cache that was \emph{least recently requested} (i.e.\ whose most recent access
is furthest in the \emph{past}). LRU is a reasonable \emph{online} heuristic because it
requires \emph{no knowledge of future requests} -- only the history of requests seen so
far -- which matches real systems where future access patterns are unknown. It is based on
the (often reasonable) assumption of \emph{temporal locality}: pages used recently are
likely to be used again soon, so pages not used recently are good eviction candidates.

\textit{Example where LRU and FIF differ:} cache size $k=2$, sequence
$0,1,2,0,1,2,0,1,2,\ldots$ (cyclic). FIF, knowing the future, can always evict the page
whose next use is furthest away (e.g.\ at each miss, keep the two pages that will be
needed soonest), achieving far fewer misses than LRU, which -- having no future
information -- ends up evicting the page that is needed \emph{immediately next} in this
adversarial cyclic pattern, causing a miss on (almost) every request. (This is exactly the
instance constructed in the Problem 18 exercise.)

\item A coin system is \textbf{canonical} if, for \emph{every} target amount $V \ge 0$,
the greedy algorithm (repeatedly subtract the largest denomination $\le$ the remaining
amount) produces a representation using the \emph{minimum possible} number of coins. The
system $\{1,5,10,25\}$ is canonical: one can check (or prove via case analysis on the
remainder modulo each denomination) that greedy's choice at each step never forces extra
coins later, for any amount. The system $\{1,3,4\}$ is \textbf{not} canonical: the
smallest counter-example is $V=6$. Greedy takes $4$, then $1$, then $1$ -- three coins
($4+1+1=6$). But the optimum is two coins: $3+3=6$.

\item \textbf{Greedy algorithm:} sort jobs in \emph{decreasing} order of profit $p_i$.
Maintain an array of time slots $1, 2, \ldots, D$ (where $D = \max_i d_i$), all initially
free. For each job $i$ in this order, find the \emph{latest} free slot $t \le d_i$ (if
any); if one exists, assign job $i$ to slot $t$ (marking it occupied) and add $p_i$ to the
total profit; otherwise, skip job $i$ (it cannot be scheduled profitably).

\textit{Why decreasing profit order is right:} we want to ``lock in'' the highest-profit
jobs first while they still have a free slot available at or before their deadline.
\textit{Why ``as late as possible'':} placing a job in the latest available slot at or
before its deadline preserves all \emph{earlier} slots for jobs with \emph{tighter
(earlier) deadlines} that may be considered later -- it never ``wastes'' an early slot on
a job that had more flexibility. Together, these ensure that whenever a later (lower-
profit) job cannot be placed, it is genuinely because all slots up to its deadline are
occupied by jobs of equal or greater profit -- which is unavoidable in any schedule.

\item All four proofs share the same high-level structure: (1) define a precise
\emph{measure of progress} or \emph{structural property} that the greedy solution
satisfies at every step (earliest finish time so far / using $\le \mathrm{depth}$
resources / no inversions and no idle time / evicting the furthest-future page); (2) show
that \emph{any} solution -- in particular, an optimal one -- can be transformed,
\emph{without making it worse}, to also satisfy this property, or show directly that
greedy's measure is always at least as good as optimal's at every step (greedy stays
ahead); (3) conclude that greedy's final objective value matches (or is no worse than) the
optimal's. In every case the proof is \emph{local} (about one step or one exchange) but
the conclusion is \emph{global} (about the entire solution), connected via induction.

\item A systematic process: (1) \textbf{Try to break it first.} Before trying to prove the
rule correct, actively search for counter-examples: construct small instances (2--4
items) designed to create \emph{ties} or \emph{near-ties} under the greedy criterion,
since these are where greedy's ``local'' view is most likely to miss a better global
choice. (2) \textbf{Look for structural invariants.} If small/random instances all agree
with brute force, try to identify \emph{what quantity} greedy is implicitly
maximizing/minimizing at each step, and check whether that quantity provably dominates the
corresponding quantity for any other solution (greedy stays ahead). (3) \textbf{Try an
exchange.} Take a hypothetical optimal solution that disagrees with greedy at the first
point of difference, and check whether swapping/adjusting it to agree with greedy at that
point can be shown to never hurt the objective. (4) \textbf{Stress test computationally.}
Generate many random small instances, compute both the greedy answer and a brute-force
optimum, and compare -- exactly as Problems 5, 9, 13, and 19 do this week. If a
discrepancy is found, minimize the counter-example by hand to understand exactly which
feature of the instance breaks the rule (often a tie, an extreme length, or a cyclic
pattern).

\end{enumerate}

\subsection*{Part B}

\begin{enumerate}[label=\textbf{B\arabic*.}, leftmargin=2.2em]

\item \textbf{False.} A defining property of greedy algorithms is that they make each
choice \emph{irrevocably} -- there is no backtracking or revision of earlier decisions.

\item \textbf{(b) Finish time, increasing.} EFT sorts by finish time $f_i$ in increasing
order; this is the key to the correctness proof (Q7--Q8).

\item Shortest-interval-first picks $[1,2]$, then $[3,4]$ (both compatible with each
other; $[0,10]$ is excluded since it overlaps both), giving a solution of size $2$. The
true optimum is also $2$ (e.g.\ $\{[1,2],[3,4]\}$, since $[0,10]$ alone only gives size
$1$) -- so on this particular instance shortest-first happens to tie the optimum. (A
larger/different instance is needed to make shortest-first \emph{strictly} worse, as
explored in Problem 6.)

\item \textbf{True.} The depth lower bound (Q10) is a property of the \emph{instance}
itself (the maximum number of pairwise-overlapping intervals at any point in time) and
does not depend on which algorithm is used -- it applies to \emph{any} valid assignment.

\item \textbf{(b) Deadline $d_i$, increasing.} EDF sorts by deadline; processing time and
profit are irrelevant to this algorithm (profit doesn't even appear in the Minimize
Lateness problem).

\item \textbf{True.} This is exactly the content of Q14 -- idle time can be removed
without increasing lateness, and inversions can be removed one at a time via an adjacent
swap that does not increase the maximum lateness.

\item The page in the cache whose \textbf{next request (use) is furthest in the future}
among all pages currently in the cache -- or any page that is never requested again, if
one is present.

\item \textbf{(b) An online algorithm that uses only the request history seen so far.}
LRU makes eviction decisions based on \emph{past} access times only, with no knowledge of
future requests -- unlike FIF, which is an offline algorithm requiring the full future
sequence.

\item $V=6$. Greedy uses \textbf{3} coins ($4+1+1$); the optimum uses \textbf{2} coins
($3+3$).

\item \textbf{True.} This greedy strategy (decreasing profit, place each job as late as
possible before its deadline) is a classic exchange-argument result and always achieves
the maximum total profit for unit-time Job Sequencing with Deadlines and Profits.

\end{enumerate}

\end{document}
